% ============================================================
%  Chapter 15
% ============================================================

\chapter{CLIO and Constraint-Leveraged Inference}

\chapterepigraph{%
  Constraints are not obstacles to thought.\\
  They are the medium in which thought becomes possible.
}{Flyxion, \textit{Semantic Infrastructure}}

\sidenote{App.\,H gives\\the formal\\admissibility\\theory}

The previous chapter developed accessibility landscapes as geometric
objects. This chapter introduces CLIO — Constraint-Leveraged Inference
Optimization — the framework for using constraint structure to make
inference tractable.

\section{Constraint as Computation}

The central claim of CLIO is counterintuitive: \emph{constraints
accelerate computation rather than impeding it}.

Without constraints, a search problem requires exploring an exponentially
large space. With tight constraints, most of this space is immediately
ruled out, and the search reduces to a small accessible region. The
tighter the constraints, the smaller the search space, and the more
efficient the inference.

\begin{definition}{Constraint-Leveraged Inference}{clio-def}
  \emph{Constraint-Leveraged Inference} (CLI) for a problem $(X, \mathcal{Q})$
  (state space $X$, query $\mathcal{Q}$) with constraint space
  $\mathcal{C} = (X, A, \kappa)$ is inference restricted to $A$:
  \[
    \mathrm{CLI}(\mathcal{Q}, \mathcal{C}) =
    \arg\max_{x \in A} P(x | \mathcal{Q}).
  \]
  The optimization is over the admissible region $A$ only.
\end{definition}

The power of CLI comes from the fact that $|A| \ll |X|$ in typical
problems. Medical diagnosis is an example: the space of all possible
symptom patterns is enormous, but the space of symptom patterns
consistent with known diseases (the admissible region) is much smaller.
A doctor's diagnostic reasoning is, informally, CLI: they do not
consider all possible diagnoses; they consider only those consistent
with the presenting symptoms.

\section{Sparse Projection}

CLIO implements constraint-leveraged inference through \emph{sparse
projection}: projecting the query into a space where most features
are zero (ruled out by constraints) and only a sparse set of features
are active.

\begin{definition}{Sparse Projection}{sparse-proj}
  A \emph{sparse projection} is a map $\pi : X \to \mathbb{R}^n$ such
  that for most queries $\mathcal{Q}$, the projected query
  $\pi(\mathcal{Q})$ has at most $k \ll n$ nonzero coordinates.
  The \emph{sparsity ratio} is $k/n$.
\end{definition}

Constraint structure naturally induces sparse projections: if the
admissible region $A$ is a low-dimensional manifold in $X$, then any
projection of $A$ into a coordinate system aligned with the constraint
directions will be sparse.

\section{Search Without Exhaustive Search}

\begin{example}[title={TCU Encoding: Constraint Restriction as Security}]
  The RSVP security duality has a precise empirical demonstration in
  analog optical neural networks. These systems use physical photonic
  circuits rather than digital processors, and suffer from hardware
  noise and low precision — which the orthodox engineering view treats
  as defects to be suppressed.

  The RSVP reinterpretation: noise and low precision are deliberate
  \emph{constraint restrictions}. By narrowing the admissible trajectory
  space (reducing $S$ locally), they eliminate the mathematical room
  that an adversarial attacker requires.

  Standard binary encoding gives attackers a \emph{gradient lever}:
  flipping the most-significant bit of an 8-bit number changes the value
  by 128, while flipping the least-significant bit changes it by 1.
  An attacker can maximize damage with minimal effort by targeting the
  highest-leverage bit.

  \emph{Truncated Complementary Unary (TCU) encoding} removes this lever
  entirely: every bit represents the same value, so no bit has more
  adversarial leverage than any other. The admissibility radius shrinks
  from $\mathcal{O}(2^{b-1})$ to $\mathcal{O}(1)$ — the mathematical
  space for a single attack collapses by an exponential factor.

  \textbf{Empirical result} on VGG-8/CIFAR-10:
  \begin{itemize}
    \item Clean baseline accuracy: $88\%$
    \item After adversarial weight attack (unprotected): $13.52\%$
    \item After RSVP-TCU defense: $86.73\%$ — near-full recovery
    \item Memory overhead of the defense: $2.36\%$ (applied to the
          $0.2\%$ most sensitive weights only)
  \end{itemize}

  This demonstrates the CLIO principle in hardware: by identifying the
  topological choke points (the high-leverage weights) and restricting
  their admissibility radius, the entire system's robustness is secured
  at minimal cost. Constraint restriction is not a limitation — it is
  precision targeted security.
\end{example}


The algorithmic consequence of sparse projection is that CLI can solve
problems that naive exhaustive search cannot.

\begin{theorem}{CLIO Complexity Bound}{clio-complexity}
  If the admissible region $A$ has volume $\mathrm{Vol}(A) = \epsilon \cdot
  \mathrm{Vol}(X)$ (a fraction $\epsilon$ of the full state space), then
  CLI inference on $A$ requires time $O(\epsilon^{-1})$ less than
  exhaustive search on $X$.
\end{theorem}

\begin{proof}[Sketch]
  Exhaustive search visits all of $X$, taking time $O(\mathrm{Vol}(X))$.
  CLI restricts to $A$, taking time $O(\mathrm{Vol}(A)) = O(\epsilon \cdot
  \mathrm{Vol}(X))$. The ratio is $O(\epsilon^{-1})$.
\end{proof}

In practice, $\epsilon$ can be astronomically small. In a protein folding
problem, the physically accessible conformations are a vanishingly small
fraction of all geometrically possible conformations. In language generation,
grammatically and semantically coherent sentences are a tiny fraction of
all possible strings. CLIO exploits this to make inference tractable.

\begin{exercise}
  Apply CLIO to the problem of finding a route between two cities on
  a road network. What is the full state space $X$? What constraints
  define the admissible region $A$ (valid routes)? Estimate the sparsity
  ratio $\epsilon$ for a road network with 1000 cities.
\end{exercise}

\begin{exercise}
  In the RSVP framework, the accessibility entropy $S(x) = \log \mathrm{Vol}(A_x)$
  measures the size of the admissible future set. Show that
  CLI inference in a state with low accessibility entropy $S$ is faster
  than in a state with high $S$. What does this say about the tradeoff
  between flexibility (high $S$) and inferential efficiency (low $S$)?
\end{exercise}

\begin{exercise}[title={Adversarial}]
  The IAD (Introspective Adversarial Distillation) framework from the
  introduction assigns a trust weight $\alpha(x) \in [0,1]$ to each
  query $x$. Reformulate this as a CLIO problem: what is the constraint
  space? What does ``admissible'' mean in the context of teacher-student
  inference? How does the sparsity ratio relate to the trust weight?
\end{exercise}

